\section{Research Objectives}

% Ý đoạn: mở phần mục tiêu nghiên cứu.
% Vai trò: chuyển từ phần tổng quan/problem sang các aims cụ thể của study.
This study is guided by the following objectives:

% Ý đoạn: liệt kê các mục tiêu nghiên cứu ở mức high-level.
% Vai trò: xác định study aims, tránh biến chương mở đầu thành checklist implementation.
\begin{enumerate}[label=\arabic*.]
    \item Develop a biologically interpretable determinant-based workflow for AMR prediction from curated genotype and phenotype resources.
    \item Investigate how different genotype-to-phenotype connection scopes affect antibiotic-specific prediction behavior.
    \item Compare direct rule-based inference with machine-learning prediction under matched biological scope settings.
    \item Evaluate the resulting prediction behavior using clinically relevant metrics, especially recall, precision, and ROC-AUC.
    \item Interpret the observed differences biologically, particularly in relation to mechanism specificity, cross-resistance, feature coverage, and determinant co-occurrence.
\end{enumerate}

% Ý đoạn: chốt lại tinh thần của bộ objectives.
% Vai trò: nhấn mạnh đây là một nghiên cứu biology-informed, không phải to-do list kỹ thuật.
Taken together, these objectives define the study as a biology-informed investigation of AMR prediction rather than a procedural modeling checklist. The emphasis is not only on predictive performance, but also on how the structure of curated genotype evidence shapes the resulting phenotype inference.

% Ý đoạn: roadmap cho phần còn lại của báo cáo.
% Vai trò: giúp người đọc biết mỗi chapter tiếp theo sẽ trả lời phần nào của bài toán.
The remainder of the report is organized as follows. Chapter 2 introduces the biological background relevant to AMR, genotype, and phenotype, Chapter 3 describes the workflow and modeling design, and Chapter 4 presents the experimental results and limitations. The report closes with a conclusion summarizing the main findings.
